\subsection{从 Fibonacci 关系呈示的尺度计数单子到自然数（商与不变量）}\label{subsec:emergent-arithmetic-fib-congruence}

本节的关键点是：\textbf{算术结构并非预置}，而是从“尺度计数的自由结构”在 Fibonacci 关系下取商（或取不变量）后被强制导出；Zeckendorf 规范形只是该商结构上的一个\textbf{可计算的规范代表系}。

令
\[
\mathcal{D}:=\Bigl\{a=(a_1,a_2,\dots): a_k\in\NN,\ \text{仅有限个 }a_k\neq 0\Bigr\}.
\]
将 $a\in\mathcal{D}$ 视为“在各尺度 $k$ 上出现 $a_k$ 次”的计数配置；$\mathcal{D}$ 在逐点加法下是自由交换幺半群（自由交换单子）。
记 $e_k\in\mathcal{D}$ 为第 $k$ 个标准基（$e_k$ 在第 $k$ 位为 $1$，其余为 $0$）。

\begin{definition}[Fibonacci 关系生成的同余]\label{def:fib-congruence}
记 $\equiv$ 为 $\mathcal{D}$ 上最小的交换幺半群同余，使得
\begin{enumerate}
  \item \textbf{（基准同余）}$e_2\equiv 2e_1$（把“同一权重的重复”折叠为下一尺度的单次出现）；
  \item \textbf{（递推同余）}对所有 $k\ge 1$ 都有
\[
e_{k+2}\ \equiv\ e_{k+1}+e_k.
\]
\end{enumerate}
等价地，$\equiv$ 由“基准折叠 + Fibonacci 递推”在“尺度计数”层面生成，并对加法封闭。
\leanverified{paper\_val\_base\_e2}
\leanverified{paper\_val\_recursion\_e}
\end{definition}

\begin{definition}[值同态（不变量）]\label{def:val-on-D}
定义映射 $\mathrm{Val}:\mathcal{D}\to\NN$ 为
\[
\mathrm{Val}(a):=\sum_{k\ge 1} a_k F_{k+1}.
\]
这是交换幺半群同态，并与前述对 $X_\infty$ 的 $\mathrm{Val}$ 定义一致（将 $c\in X_\infty$ 视为 $a_k=c_k$ 的特殊情形）。
\leanpartial{valE}{仅对单基底向量 e\_k 形式化；一般 $\mathcal{D}$ 元素上的 Val 尚未形式化}
\end{definition}

\begin{proposition}[Fibonacci 同余的完备性]\label{prop:val-invariant}
若 $a\equiv b$，则 $\mathrm{Val}(a)=\mathrm{Val}(b)$。因此 $\mathrm{Val}$ 在商幺半群 $\mathcal{D}/\!\equiv$ 上诱导良定义的同态
\[
\overline{\mathrm{Val}}:\mathcal{D}/\!\equiv\ \to\ \NN.
\]
\leanverified{paper\_zeckendorf\_val\_invariant}
\end{proposition}

\textbf{证明：}
只需验证生成关系保持值。对递推同余 $e_{k+2}\equiv e_{k+1}+e_k$，
\[
\mathrm{Val}(e_{k+2})=F_{k+3}=F_{k+2}+F_{k+1}=\mathrm{Val}(e_{k+1}+e_k),
\]
对基准同余 $e_2\equiv 2e_1$，
\[
\mathrm{Val}(e_2)=F_3=2=2F_2=\mathrm{Val}(2e_1).
\]
其余由同余对加法封闭、以及 $\mathrm{Val}$ 的加法同态性推出。

\begin{theorem}[自然数的涌现：商幺半群的同构]\label{thm:monoid-quotient-is-N}
映射 $\overline{\mathrm{Val}}:\mathcal{D}/\!\equiv \to \NN$ 是交换幺半群同构。并且 $X_\infty$（Zeckendorf 合法序列）给出每个同余类的唯一规范代表：对任意 $a\in\mathcal{D}$，存在唯一 $c\in X_\infty$ 使得 $a\equiv c$ 且 $\mathrm{Val}(a)=\mathrm{Val}(c)$。
\leanpartial{paper\_val\_e\_index\_eq\_fib\_succ}{已形式化标定引理 $[e\_k]=F\_{k+1}\varepsilon$ 的值侧；同构与规范代表唯一性尚未形式化}
\leanverified{valE\_zero\_one\_two\_table}
\leanverified{valE\_pos\_of\_ge\_one}
\leanverified{valE\_mono}
\leanverified{valE\_strict\_mono\_of\_ge\_two}
\leanverified{paper\_valE\_milestone\_package}
\leanverified{paper\_monoid\_quotient\_is\_n}
\end{theorem}

\textbf{证明：}
\begin{itemize}
  \item \textbf{满射性：}对任意 $N\in\NN$，取其 Zeckendorf 数字 $Z(N)=(c_1,c_2,\dots)\in X_\infty$（\cite{Zeckendorf1972}），则 $\mathrm{Val}(Z(N))=N$，故 $\overline{\mathrm{Val}}$ 满射。
  \item \textbf{同构性（用呈示把一切类压到单生成元）：}
  在商幺半群 $\mathcal{D}/\!\equiv$ 中，令 $\varepsilon:=[e_1]$。由基准同余有 $[e_2]=2\varepsilon$。再由递推同余归纳得对所有 $k\ge 1$，
  \[
  [e_{k+2}]=[e_{k+1}]+[e_k],
  \qquad
  \Rightarrow\qquad
  [e_k]=F_{k+1}\,\varepsilon.
  \]
  因此对任意 $a=\sum_{k\ge 1} a_k e_k\in\mathcal{D}$，
  \[
  [a]=\sum_{k\ge 1} a_k [e_k]=\left(\sum_{k\ge 1} a_k F_{k+1}\right)\varepsilon
  =\mathrm{Val}(a)\,\varepsilon.
  \]
这说明每个同余类都由其值 $\mathrm{Val}(a)$ 完全刻画，且 $\overline{\mathrm{Val}}([a])=\mathrm{Val}(a)$ 正是把 $n\varepsilon$ 送到 $n$ 的单调同构。于是 $\overline{\mathrm{Val}}$ 为交换幺半群同构。
  \item \textbf{规范代表：}对任意 $a\in\mathcal{D}$，令 $N=\mathrm{Val}(a)$，取 $c=Z(N)\in X_\infty$。上一步给出
  \(
  [a]=N\varepsilon=[c],
  \)
  即 $a\equiv c$。唯一性由 Zeckendorf 唯一性推出：若 $c,c'\in X_\infty$ 且 $c\equiv c'$，则 $\mathrm{Val}(c)=\mathrm{Val}(c')$，从而 $c=c'$。
\end{itemize}

\begin{corollary}[折叠作为商映射的可计算截面]\label{cor:fold-as-section}
上一定理的结论可概括为：
\[
\mathcal{D}\ \twoheadrightarrow\ \mathcal{D}/\!\equiv\ \cong\ \NN
\quad\text{且}\quad
X_\infty \hookrightarrow \mathcal{D}
\]
是一个可计算的截面（选择规范代表）。这正是”折叠（folding）= 取不变量/取商 + 选规范形”的抽象机制在 Fibonacci 情形下的具体实现。
\leanverified{paper\_ea\_fold\_as\_section}
\end{corollary}
\begin{theorem}[有限观察逆系统的折叠下降与坏例预算]\label{distill:connes_consani_arithmetic_site-emergent_arithmetic_from_action_site_sections-finite-observer-descent-budget}
对每个整数 \(q\ge 2\)，定义有限分类观察截面
\[
\rho_q:\mathcal{D}\longrightarrow \mathbb{Z}/q\mathbb{Z},
\qquad
\rho_q(a):=\left[\mathrm{Val}(a)\right]_q
=\left[\sum_{k\ge 1}a_kF_{k+1}\right]_q .
\]
当 \(q\mid q'\) 时，记 \(r_{q',q}:\mathbb{Z}/q'\mathbb{Z}\to \mathbb{Z}/q\mathbb{Z}\) 为自然约化映射。令
\[
L_M:=\operatorname{lcm}(2,3,\dots,M)\qquad (M\ge 2),
\]
并称 \(a,b\in\mathcal{D}\) 为 \(M\)-不可分辨，如果 \(\rho_q(a)=\rho_q(b)\) 对所有 \(2\le q\le M\) 成立。则以下结论成立。
\begin{enumerate}
  \item 每个 \(\rho_q\) 都是 \(\mathcal{D}\) 的平移作用上的加法截面读数：对任意 \(a,c\in\mathcal{D}\)，
  \[
  \rho_q(a+c)=\rho_q(a)+\rho_q(c).
  \]
  且 \(\rho_q\) 与 Fibonacci 折叠关系相容：
  \[
  \rho_q(e_2)=\rho_q(2e_1),
  \qquad
  \rho_q(e_{k+2})=\rho_q(e_{k+1}+e_k)\quad(k\ge 1).
  \]
  因此 \(\rho_q\) 下降为商上的截面
  \[
  \overline{\rho}_q:\mathcal{D}/\!\equiv\longrightarrow \mathbb{Z}/q\mathbb{Z},
  \]
  并且这些截面对观察覆盖的细化构成稳定逆系统：若 \(q\mid q'\)，则
  \[
  r_{q',q}\circ \overline{\rho}_{q'}=\overline{\rho}_q.
  \]
  \item 对任意 \(a,b\in\mathcal{D}\)，\(a,b\) 为 \(M\)-不可分辨，当且仅当
  \[
  \mathrm{Val}(a)-\mathrm{Val}(b)\in L_M\mathbb{Z}.
  \]
  特别地，若 \(a,b\) 为 \(M\)-不可分辨且 \(a\not\equiv b\)，则
  \[
  \bigl|\mathrm{Val}(a)-\mathrm{Val}(b)\bigr|\ge L_M.
  \]
  \item 设 \(B<L_M\)，并令
  \[
  \mathcal{D}_{\le B}:=\{a\in\mathcal{D}:\mathrm{Val}(a)\le B\}.
  \]
  则有限观察覆盖 \(\{\rho_q:2\le q\le M\}\) 在 \(\mathcal{D}_{\le B}\) 上已经充分：若 \(a,b\in\mathcal{D}_{\le B}\) 且 \(a,b\) 为 \(M\)-不可分辨，则 \(a\equiv b\)。
  \item 在同一预算假设 \(B<L_M\) 下，令
  \[
  \mathsf{Add}_q(a,b):=\rho_q(a)+\rho_q(b),
  \qquad
  \mathsf{Mul}_q(a,b):=\rho_q(a)\rho_q(b)
  \]
  为 \(\mathbb{Z}/q\mathbb{Z}\) 中的局部加法与乘法读数。若 \(\mathrm{Val}(a)+\mathrm{Val}(b)\le B\)，且 \(u\in\mathcal{D}_{\le B}\) 满足
  \[
  \rho_q(u)=\mathsf{Add}_q(a,b)\qquad(2\le q\le M),
  \]
  则 \(u\equiv a+b\)。若 \(\mathrm{Val}(a)\mathrm{Val}(b)\le B\)，且 \(u\in\mathcal{D}_{\le B}\) 满足
  \[
  \rho_q(u)=\mathsf{Mul}_q(a,b)\qquad(2\le q\le M),
  \]
  则
  \[
  u\equiv \bigl(\mathrm{Val}(a)\mathrm{Val}(b)\bigr)e_1.
  \]
  因而，在固定稳定值预算内，加法、乘法与同余判定均由平移作用、Fibonacci 折叠相容性、有限分类观察截面及其逆系统细化唯一重构。
\end{enumerate}
\end{theorem}
\begin{theorem}[max-plus 包络下的有限相位下降与反例骨架]\label{distill:connes_consani_arithmetic_site-characteristic_one_stable_value_reconstruction-max-plus-envelope-phase-descent}
令
\(
\mathbb T_{\max}:=\RR\cup\{-\infty\}
\)
带有运算
\(
x\oplus y:=\max\{x,y\}
\)
与
\(
x\odot y:=x+y
\)，并约定 \(2^{-\infty}=0\)。
对 \(a\in\mathcal D\)，记
\[
\operatorname{supp}(a):=\{k\ge1:a_k\ne0\},
\qquad
s(a):=\#\operatorname{supp}(a),
\]
并定义其 Fibonacci 尺度的 max-plus 包络为
\[
\mu(a):=\bigoplus_{k\in\operatorname{supp}(a)}
\bigl(\log_2 a_k\odot\log_2 F_{k+1}\bigr)
=\max_{k\in\operatorname{supp}(a)}\log_2(a_kF_{k+1}),
\]
其中空最大值取为 \(-\infty\)。
给定 \(M\ge2\)，令
\[
L_M:=\operatorname{lcm}(2,3,\ldots,M),
\qquad
\rho_q(a):=[\operatorname{Val}(a)]_q\quad(2\le q\le M).
\]
若 \(a,b\in\mathcal D\) 满足
\[
\rho_q(a)=\rho_q(b)\qquad(2\le q\le M)
\]
以及尺度--相位转移不等式
\[
\max\{s(a),s(b)\}\,2^{\max\{\mu(a),\mu(b)\}}<L_M,
\]
则
\[
a\equiv b.
\]
因此，任意满足所有有限相位读数相同但 \(a\not\equiv b\) 的反例必满足
\[
\max\{s(a),s(b)\}\,2^{\max\{\mu(a),\mu(b)\}}\ge L_M.
\]
该下界是尖锐的：对每个 \(M\ge2\)，取
\(
a=L_Me_1
\)
与
\(
b=0
\)，则 \(\rho_q(a)=\rho_q(b)\) 对所有 \(2\le q\le M\) 成立，且 \(a\not\equiv b\)，并且
\[
\max\{s(a),s(b)\}\,2^{\max\{\mu(a),\mu(b)\}}=L_M.
\]
进一步，令 \(I\subset\NN_{\ge1}\) 为有限集，\(B\in\NN\)，并令
\[
\mathcal D_{I,B}:=\{a\in\mathcal D:\operatorname{supp}(a)\subset I,\ 0\le a_k\le 2^B\ \text{对所有 }k\in I\}.
\]
对 \(a\in\mathcal D_{I,B}\) 定义有限 max-plus 影子
\[
\Theta_{I,B}(a)_k:=
\begin{cases}
-\infty,&a_k=0,\\
\lceil\log_2 a_k\rceil,&a_k>0,
\end{cases}
\qquad(k\in I).
\]
若序列 \((a^{(t)})_{t\ge0}\subset\mathcal D_{I,B}\) 的 \(\Theta_{I,B}(a^{(t)})\) 在 \(\mathbb T_{\max}^{I}\) 的逐坐标序中单调不降，则该影子至多经过 \(|I|(B+1)\) 次严格变化后稳定。若此外每个 \(\rho_q(a^{(t)})\) 都最终稳定，且
\[
|I|\,2^B\max_{k\in I}F_{k+1}<L_M,
\]
则商类 \([a^{(t)}]\in\mathcal D/\!\equiv\) 最终稳定。
\end{theorem}
\begin{proof}
设
\[
C:=\max\{s(a),s(b)\}\,2^{\max\{\mu(a),\mu(b)\}}.
\]
由 \(\mu\) 的定义，对每个 \(k\in\operatorname{supp}(a)\) 有
\(
a_kF_{k+1}\le 2^{\mu(a)}\le 2^{\max\{\mu(a),\mu(b)\}}
\)。
因此
\[
0\le \operatorname{Val}(a)
=\sum_{k\ge1}a_kF_{k+1}
\le s(a)2^{\mu(a)}\le C.
\]
同理 \(0\le\operatorname{Val}(b)\le C\)。若 \(C<L_M\)，则
\[
|\operatorname{Val}(a)-\operatorname{Val}(b)|<L_M.
\]
另一方面，\(\rho_q(a)=\rho_q(b)\) 对所有 \(2\le q\le M\) 成立，故整数
\(
\operatorname{Val}(a)-\operatorname{Val}(b)
\)
被每个 \(q=2,\ldots,M\) 整除，因而被 \(L_M\) 整除。结合上式只能得到
\[
\operatorname{Val}(a)=\operatorname{Val}(b).
\]
由定理 \ref{thm:monoid-quotient-is-N}，\(\overline{\operatorname{Val}}:\mathcal D/\!\equiv\to\NN\) 为同构，故 \([a]=[b]\)，即 \(a\equiv b\)。反例下界即为该结论的逆否命题。

对尖锐性，若 \(a=L_Me_1\) 且 \(b=0\)，则 \(\operatorname{Val}(a)=L_MF_2=L_M\)，而 \(\operatorname{Val}(b)=0\)。由于每个 \(q\in\{2,\ldots,M\}\) 都整除 \(L_M\)，有 \(\rho_q(a)=\rho_q(b)\)。但 \(\operatorname{Val}(a)\ne\operatorname{Val}(b)\)，再由定理 \ref{thm:monoid-quotient-is-N} 得 \(a\not\equiv b\)。同时 \(s(a)=1\)、\(s(b)=0\)、\(F_2=1\)，所以
\[
\mu(a)=\log_2(L_M),
\qquad
\max\{s(a),s(b)\}\,2^{\max\{\mu(a),\mu(b)\}}=L_M.
\]

最后考虑 \((a^{(t)})_{t\ge0}\subset\mathcal D_{I,B}\)。每个坐标 \(\Theta_{I,B}(a^{(t)})_k\) 只取有限链
\(
\{-\infty,0,1,\ldots,B\}
\)
中的值。若 \(\Theta_{I,B}(a^{(t)})\) 逐坐标单调不降，则每个坐标最多严格增加 \(B+1\) 次，故整个影子最多严格变化 \(|I|(B+1)\) 次，因而最终稳定。

设所有 \(\rho_q(a^{(t)})\) 在 \(t\ge T_0\) 后稳定。对任意 \(t,t'\ge T_0\)，有
\(
\rho_q(a^{(t)})=\rho_q(a^{(t')})
\)
对所有 \(2\le q\le M\) 成立。又因 \(a^{(t)},a^{(t')}\in\mathcal D_{I,B}\)，有
\[
s(a^{(t)})\le |I|,
\qquad
2^{\mu(a^{(t)})}\le 2^B\max_{k\in I}F_{k+1},
\]
并且 \(a^{(t')}\) 同样满足该界。由假设
\(
|I|2^B\max_{k\in I}F_{k+1}<L_M
\)
及已经证明的下降判据可得
\(
a^{(t)}\equiv a^{(t')}
\)。因此 \([a^{(t)}]\) 在 \(t\ge T_0\) 后恒定。证毕。
\end{proof}

\begin{proof}
首先，\(\mathrm{Val}\) 是交换幺半群同态，故对任意 \(a,c\in\mathcal{D}\) 有
\[
\rho_q(a+c)=[\mathrm{Val}(a+c)]_q
=[\mathrm{Val}(a)+\mathrm{Val}(c)]_q
=\rho_q(a)+\rho_q(c).
\]
这正是 \(\mathcal{D}\) 对自身平移作用下的截面相容性。对折叠生成元，利用本节 Fibonacci 标定，得到
\[
\rho_q(e_2)=[F_3]_q=[2F_2]_q=\rho_q(2e_1),
\]
且对所有 \(k\ge 1\)，
\[
\rho_q(e_{k+2})=[F_{k+3}]_q=[F_{k+2}+F_{k+1}]_q
=\rho_q(e_{k+1}+e_k).
\]
因此 \(\rho_q\) 对生成同余关系取相同值，并由同余对加法封闭可知 \(\rho_q\) 在 \(\mathcal{D}/\!\equiv\) 上诱导良定义映射 \(\overline{\rho}_q\)。若 \(q\mid q'\)，则对任意 \(a\in\mathcal{D}\)，
\[
(r_{q',q}\circ\rho_{q'})(a)=r_{q',q}([\mathrm{Val}(a)]_{q'})=[\mathrm{Val}(a)]_q=\rho_q(a),
\]
所以下降后仍有 \(r_{q',q}\circ\overline{\rho}_{q'}=\overline{\rho}_q\)。这证明第一项。

接着，\(a,b\) 为 \(M\)-不可分辨，当且仅当对每个 \(2\le q\le M\)，
\[
[\mathrm{Val}(a)-\mathrm{Val}(b)]_q=0,
\]
也即每个 \(q\in\{2,\dots,M\}\) 都整除整数 \(\mathrm{Val}(a)-\mathrm{Val}(b)\)。这等价于其最小公倍数 \(L_M\) 整除该差值，即
\[
\mathrm{Val}(a)-\mathrm{Val}(b)\in L_M\mathbb{Z}.
\]
若同时 \(a\not\equiv b\)，由定理 \ref{thm:monoid-quotient-is-N}，商类由 \(\mathrm{Val}\) 完全刻画，故 \(\mathrm{Val}(a)\ne \mathrm{Val}(b)\)。于是该非零差值是 \(L_M\) 的整数倍，从而其绝对值至少为 \(L_M\)。第二项得证。

若 \(a,b\in\mathcal{D}_{\le B}\) 且 \(B<L_M\)，则
\[
\bigl|\mathrm{Val}(a)-\mathrm{Val}(b)\bigr|\le B<L_M.
\]
若 \(a,b\) 为 \(M\)-不可分辨，第二项给出 \(L_M\mid \mathrm{Val}(a)-\mathrm{Val}(b)\)，因此只能有 \(\mathrm{Val}(a)=\mathrm{Val}(b)\)。再由定理 \ref{thm:monoid-quotient-is-N}，得到 \(a\equiv b\)。第三项得证。

最后证明局部运算读数的唯一性。若 \(\mathrm{Val}(a)+\mathrm{Val}(b)\le B\) 且 \(u\in\mathcal{D}_{\le B}\) 满足 \(\rho_q(u)=\mathsf{Add}_q(a,b)\) 对所有 \(2\le q\le M\) 成立，则由 \(\rho_q\) 的加法性，
\[
\rho_q(u)=\rho_q(a)+\rho_q(b)=\rho_q(a+b).
\]
同时 \(a+b\in\mathcal{D}_{\le B}\)。由第三项应用于 \(u\) 与 \(a+b\)，得 \(u\equiv a+b\)。

若 \(\mathrm{Val}(a)\mathrm{Val}(b)\le B\)，令
\[
v:=\bigl(\mathrm{Val}(a)\mathrm{Val}(b)\bigr)e_1\in\mathcal{D}.
\]
由于本节标定下 \(F_2=1\)，有
\[
\mathrm{Val}(v)=\mathrm{Val}(a)\mathrm{Val}(b)\le B.
\]
并且对每个 \(q\)，
\[
\rho_q(v)=[\mathrm{Val}(a)\mathrm{Val}(b)]_q
=[\mathrm{Val}(a)]_q[\mathrm{Val}(b)]_q
=\rho_q(a)\rho_q(b)=\mathsf{Mul}_q(a,b).
\]
若 \(u\in\mathcal{D}_{\le B}\) 具有相同的乘法局部读数，则 \(u\) 与 \(v\) 在所有 \(2\le q\le M\) 的观察截面上不可分辨。第三项给出 \(u\equiv v\)，即
\[
u\equiv \bigl(\mathrm{Val}(a)\mathrm{Val}(b)\bigr)e_1.
\]
同余判定本身即第三项。第四项随之成立。\end{proof}

\textbf{证明：}
由定理 \ref{thm:monoid-quotient-is-N} 的同构性与唯一规范代表结论直接推出。
